\documentclass[11pt, a4paper]{article}

\usepackage{fontspec}
\usepackage[spanish, bidi=basic, provide=*]{babel}
\usepackage{amsmath}
\usepackage{amssymb}
\usepackage{graphicx}
\usepackage{booktabs}
\usepackage{microtype}
\usepackage{enumitem}
\usepackage[style=numeric, backend=biber]{biblatex}
\usepackage{hyperref}
\usepackage{cleveref}

\addbibresource{./references.bib}

\title{Metodología para la Construcción de una Celda Solar Sensibilizada por
Colorante utilizando Materiales Accesibles}

\author{Juan Acuña \and Julian Avila \and Camilo Huertas
\and Bryan Martinez \and Adriano Parada \and David Rodriguez}
\date{}

\begin{document}

\maketitle

\section{Introducción}

Este documento expone la metodología para sintetizar y ensamblar una celda
solar sensibilizada por colorante. El dispositivo opera mediante la inyección
de carga fotoinducida desde un sensibilizador molecular orgánico hacia la banda
de conducción de un semiconductor de brecha ancha. El protocolo emplea
materiales accesibles que sustituyen reactivos de grado analítico e
instrumentos especializados, preservando rigurosamente los procesos
fisicoquímicos fundamentales.

\section{Objetivos}

El objetivo general es construir y caracterizar una celda solar fotovoltaica
funcional mediante materiales accesibles, ilustrando los principios de
conversión de energía solar y la dinámica de transferencia electrónica en
interfaces heterogéneas. Los objetivos específicos comprenden la formulación
de una suspensión coloidal homogénea de nanopartículas de dióxido de titanio
con un tensoactivo comercial para formar una película delgada uniforme; la
inducción de sinterización en la película semiconductora mediante tratamiento
térmico para establecer conectividad eléctrica; la extracción de clorofila de
tejido vegetal como sensibilizador óptico primario; y el ensamblaje de la
celda electroquímica para evaluar sus parámetros eléctricos fundamentales bajo
radiación lumínica.

\section{Materiales}

El experimento requiere sustratos conductores de vidrio y polvo de dióxido de
titanio en fase anatasa. Se emplean adicionalmente jabón líquido transparente
como agente tensoactivo, solución de ácido acético como medio de dispersión,
hojas frescas de espinaca como fuente de clorofila y alcohol etílico como
solvente orgánico, junto con una fuente térmica, grafito, clips metálicos de
presión y un electrolito de yodo-yoduro que actúa como mediador de
transferencia de carga.

\section{Metodología de Ensamblaje}

La construcción de la celda se ejecuta en etapas secuenciales que incluyen la
formación de la matriz semiconductora, la funcionalización óptica, la
preparación del electrodo catalítico y el cierre del circuito electroquímico.

\subsection{Día 1}
\subsubsection{Preparación y Deposición del Semiconductor}

El proceso inicia sintetizando la matriz mesoporosa para el transporte
electrónico. Se dispersa polvo de dióxido de titanio en ácido acético dentro
de un mortero. La adición de jabón líquido introduce moléculas anfifílicas que
reducen la tensión superficial y proveen estabilización estérica, previniendo
la aglomeración y generando una suspensión coloidal isotrópica. Se identifica
la cara conductora del sustrato de vidrio mediante multímetro y se delimita el
área activa con cinta adhesiva. Tras depositar la suspensión y aplicar el
tratamiento térmico, el sustrato se enfría gradualmente para evitar fracturas
por estrés térmico.

\subsubsection{Extracción de Clorofila a partir de Espinaca}

La absorción de fotones exige la adsorción de un cromóforo sobre el
semiconductor. Las hojas de espinaca (\textit{Spinacia oleracea}) proveen los
complejos antena pigmento-proteína, de las cuales se descarta la nervadura
central al carecer de pigmentos. Se macera el tejido residual con alcohol
etílico para incrementar la relación área-volumen. La acción mecánica y el
solvente polar rompen las membranas tilacoidales, liberando la clorofila. El
etanol solubiliza los anillos de porfirina, y la decantación posterior separa
los residuos celulares insolubles, aislando una solución homogénea lista para
la sensibilización.

\subsubsection{Sinterizado Térmico}

Se emplea una llama azul de alta temperatura producida por un mechero o lámpara
de alcohol; una llama amarilla genera hollín que contamina el dióxido de
titanio de forma irreversible. El sustrato se ubica a cuatro centímetros sobre
la llama mediante un soporte universal para recibir calor radiante. Este
calentamiento se mantiene durante quince minutos; la carbonización inicial del
tensoactivo oscurece la película, la cual recobra un color blanco puro al
oxidarse el carbono, indicando la fusión y sinterización de las nanopartículas
en una red tridimensional interconectada. Finalmente, el enfriamiento debe
ocurrir gradualmente. El enfriamiento por convección ambiental dura al menos
quince minutos; una remoción abrupta del calor induce choque térmico y
delaminación del semiconductor.

\subsection{Día 2}
\subsubsection{Sensibilización del Semiconductor}

El electrodo de dióxido de titanio a temperatura ambiente se sumerge en la
solución de clorofila. Los grupos funcionales del pigmento interactúan con los
sitios de coordinación insatisfechos del semiconductor. Esta quimisorción
asegura un fuerte acoplamiento electrónico entre el estado excitado del
pigmento y la banda de conducción. La absorción óptica induce una transición
electrónica e inyecta carga eficientemente hacia los estados desocupados del
óxido.

\subsubsection{Preparación del Contraelectrodo}

El circuito requiere un electrodo catalítico para reducir el mediador redox. Se
identifica la cara conductora de un segundo sustrato de vidrio y se deposita
una capa de carbono. Este depósito se logra mediante fricción con grafito o
exposición a hollín. Este catalizador disminuye el sobrepotencial necesario
para la transferencia electrónica hacia el electrolito.

\subsubsection{Ensamblaje y Evaluación del Dispositivo}

El dispositivo final adopta una configuración de capas apiladas. El electrodo
fotosensibilizado y el contraelectrodo se superponen con sus caras activas en
contacto directo. Se aplica un desplazamiento lateral para exponer las
terminales conductoras externas. El sistema se estabiliza mecánicamente con
clips metálicos de presión. Se instila electrolito de yodo-yoduro en el
espacio intermedio. El líquido se distribuye por capilaridad a través de la
matriz porosa para completar el circuito iónico. La caracterización final bajo
radiación fotónica mediante multímetro debe registrar una diferencia de
potencial asimétrica y una corriente de cortocircuito, demostrando la
extracción de trabajo eléctrico del campo de radiación incidente.

\printbibliography
\nocite{*}
\end{document}
